\section{Collatz prime-flow branch: exact cocycle, prime flux, and holonomy}

\subsection{Shortcut Collatz map}

Define
\begin{equation}
C(n)=
\begin{cases}
n/2,&n\equiv0\pmod2,\\[2mm]
(3n+1)/2,&n\equiv1\pmod2,
\end{cases}
\qquad n\in\mathbb N_{>0}.
\end{equation}
For a valid step \(n_k\to n_{k+1}=C(n_k)\), define
\begin{equation}
r_k=\frac{n_{k+1}}{n_k},
\qquad
b_k=n_k\bmod2\in\{0,1\}.
\end{equation}
The implementation stores \(r_k\) first as an exact reduced rational.

\subsection{The \(\mathbb Z_2\) branch representation}

The non-trivial one-dimensional unitary representation of \(\mathbb Z_2\) is
\begin{equation}
\boxed{\rho(b)=(-1)^b=e^{i\pi b}.}
\end{equation}
Define the Mellin-sign step cocycle
\begin{equation}
\boxed{w_s(n_k)=r_k^{-s}(-1)^{b_k}.}
\label{eq:step-cocycle}
\end{equation}
For \(s=\sigma+i\tau\),
\begin{equation}
|w_s|=r_k^{-\sigma},
\qquad
\arg w_s=-\tau\log r_k+\pi b_k\pmod{2\pi}.
\end{equation}
Status: \status{PROJECT\_EXACT} after definition.

\subsection{Orbit factorisation theorem}

Let
\begin{equation}
\mathcal O_m=(n_0,n_1,\ldots,n_m),
\qquad n_{k+1}=C(n_k),
\end{equation}
and
\begin{equation}
W_s(\mathcal O_m)=\prod_{k=0}^{m-1}w_s(n_k),
\qquad
N_{\rm odd}=\sum_{k=0}^{m-1}b_k.
\end{equation}
Then
\begin{equation}
\boxed{
W_s(\mathcal O_m)
=
\left(\frac{n_m}{n_0}\right)^{-s}
(-1)^{N_{\rm odd}}.
}
\label{eq:orbit-factor}
\end{equation}
Status: \status{PROJECT\_EXACT}. The proof is the telescoping product
\begin{equation}
\prod_{k<m}\frac{n_{k+1}}{n_k}=\frac{n_m}{n_0}
\end{equation}
combined with multiplication of the parity signs.

\subsection{Complementary defect and the half-axis}

Set
\begin{equation}
L=\log\frac{n_m}{n_0}.
\end{equation}
From \eqref{eq:orbit-factor},
\begin{equation}
\log|W_s|=-\sigma L,
\qquad
\log|W_{1-s}|=-(1-\sigma)L.
\end{equation}
Define
\begin{equation}
\boxed{
D_C(s;\mathcal O_m)
=
\big(\log|W_s|-\log|W_{1-s}|\big)^2
=(2\sigma-1)^2L^2.
}
\label{eq:collatz-defect}
\end{equation}
For a non-closed path \(L\neq0\),
\begin{equation}
\boxed{D_C=0\iff \RePart(s)=\frac12.}
\end{equation}
For a closed path \(L=0\), \(D_C=0\) for every \(s\). Closed paths are therefore degenerate controls and cannot serve as evidence that the half-axis is uniquely selected.

\subsection{Critical-line conjugate pair}

For \(s=\tfrac12+i\tau\),
\begin{equation}
\boxed{W_{1-s}=\overline{W_s}.}
\end{equation}
This is the exact origin of an 18-pair \(\Tthirtysix\) construction used below.

\subsection{Prime-valuation flux}

Define the stepwise prime flux
\begin{equation}
\boxed{\Delta v_{p,k}=v_p(n_{k+1})-v_p(n_k).}
\label{eq:prime-flux}
\end{equation}
By unique factorisation,
\begin{equation}
\boxed{r_k=\prod_p p^{\Delta v_{p,k}},}
\qquad
\log r_k=\sum_p\Delta v_{p,k}\log p.
\end{equation}
Hence the exact step cocycle factorises prime by prime:
\begin{equation}
\boxed{
w_s(n_k)
=(-1)^{b_k}\prod_p p^{-s\Delta v_{p,k}}.
}
\label{eq:prime-cocycle}
\end{equation}
For the full orbit define
\begin{equation}
\Delta V_p=\sum_{k<m}\Delta v_{p,k}=v_p(n_m)-v_p(n_0).
\end{equation}
Then
\begin{equation}
\boxed{
W_s(\mathcal O_m)=(-1)^{N_{\rm odd}}\prod_p p^{-s\Delta V_p}.
}
\label{eq:prime-orbit}
\end{equation}
Only finitely many \(\Delta V_p\) are non-zero for a finite endpoint pair.

On the critical line, the orbit phase becomes
\begin{equation}
\boxed{
\Theta_{\mathcal O}(\tau)
=
\pi N_{\rm odd}
-	au\sum_p\Delta V_p\log p
\pmod{2\pi}.
}
\label{eq:prime-phase}
\end{equation}
Thus
\begin{equation}
\Theta_{\mathcal O}=\Theta_{\mathbb Z_2}+\Theta_{\rm prime}.
\end{equation}

\subsection{Closed-cycle cancellation of prime flux}

If \(n_m=n_0\), then \(\Delta V_p=0\) for every prime, so
\begin{equation}
\Theta_{\rm prime}=0,
\qquad
\boxed{W_s=(-1)^{N_{\rm odd}}.}
\end{equation}
For the control cycle
\begin{equation}
1\to2\to1,
\end{equation}
there is one odd source state, so
\begin{equation}
\boxed{W_s=-1=e^{i\pi}.}
\end{equation}
Two traversals give \(+1\). The one-loop/double-loop closure is therefore exact for the defined sign cocycle.

\subsection{Where the non-trivial prime interaction resides}

The even branch \(n\mapsto n/2\) is multiplicative and transparent in valuation coordinates. The odd branch contains
\begin{equation}
3n+1,
\end{equation}
and addition prevents factorwise separation:
\begin{equation}
v_p(3n+1)\neq v_p(3)+v_p(n)
\end{equation}
in general. Thus the map
\begin{equation}
\mathbf v(n)\mapsto\mathbf v(3n+1)-\mathbf e_2
\end{equation}
is an arithmetic coupling of prime channels.

\INTERP
Zeta's Euler product represents free multiplicative composition of prime powers in its convergence region, whereas the Collatz odd branch deterministically re-couples prime valuations through the additive operation \(3n+1\). This is a useful physical-style analogy (``free modes'' versus ``interaction vertex''), not a theorem identifying a quantum field theory with Collatz arithmetic.

\subsection{Weighted Koopman operator}

Define, on a declared function space of observables,
\begin{equation}
\boxed{(U_s f)(n)=w_s(n)f(C(n)).}
\end{equation}
Then
\begin{equation}
U_s^m f(n_0)=W_s(\mathcal O_m)f(n_m).
\end{equation}
Operator approaches to the \(3x+1\) system, including Koopman methods and Fourier analysis of sign sequences, exist independently in the literature \cite{LeventidesPoulios2020}. The specific Mellin-sign cocycle and prime-flux PhaseNav embedding in this monograph are project constructions.
